\sectionCenteredToc{Заключение}
\label{sec:outro}

В данном курсовом проекте было разработано программное средство для управления файловой системой, которое позволяет эффективно навигироваться по локальным и сетевым ресурсам, просматривать содержимое директорий и управлять файлами. В ходе работы над проектом были изучены основные аспекты разработки программного обеспечения для платформы Windows, включая анализ требований, проектирование модульной архитектуры, реализацию с использованием Win32 API и тестирование.

Приложение создано с использованием современных технологий и инструментов разработки на языке C++, что позволяет обеспечить его высокую производительность и надежность. При разработке были реализованы следующие основные функции программного средства:

\begin{itemize}
    \item модуль для работы с файловой системой;
    \item модуль для работы с сетевыми устройствами;
    \item модуль управления настройками приложения;
    \item модуль пользовательского интерфейса;
    \item модуль утилитарных функций;
    \item алгоритм загрузки настроек из реестра Windows;
    \item алгоритм расширения узлов дерева каталогов;
    \item алгоритм обновления строки состояния директории;
    \item алгоритм обработки файлов директории;
    \item алгоритм добавления корневых дисков в дерево;
    \item алгоритм поиска сетевых устройств;
    \item алгоритм получения списка сетевых шаров;
\end{itemize}

Для достижения поставленных целей были использованы современные подходы и методологии разработки программного обеспечения, такие как модульная архитектура, принцип единственной ответственности (Single Responsibility Principle) и инкапсуляция данных. Это позволило обеспечить высокое качество программного продукта, его поддерживаемость и расширяемость.

В ходе работы над проектом были выявлены некоторые проблемы и ограничения, которые могут быть устранены в будущем. Например, добавление функциональности копирования и перемещения файлов, реализация поиска файлов по имени и содержимому, улучшение обработки больших директорий с тысячами файлов, оптимизация работы с сетевыми ресурсами при медленном соединении.